\section{Цель работы}
Приобретение практических навыков создания процессов операционной системы и организации обмена данными между ними посредством неименованных каналов.

\section{Задание}
Разработать две отдельные программы: родительский и дочерний процессы. Первая строка, введённая пользователем в консоль родительского процесса, содержит имя выходного файла. Родитель передаёт это имя дочернему процессу при его запуске.

Последующие команды имеют вид <<число и перевод строки>>. Родитель передаёт команды через первый канал, подключённый к стандартному входу дочернего процесса. Дочерний процесс проверяет число на простоту. Составные числа записываются в выходной файл. При получении отрицательного или простого числа оба процесса завершают работу. Проверка числа выполняется на стороне дочернего процесса.

Второй канал используется для передачи статусов от дочернего процесса к родительскому. Необходимо обрабатывать ошибки системных операций и освобождать используемые ресурсы~\cite{task}.

\subsection{Принятые уточнения}
Реализация выполнена на C++17 с собственным платформенным интерфейсом над API операционных систем. Стандартная библиотека применяется для хранения данных, разбора строк и работы с консолью; создание процессов, каналов и запись выходного файла выполняются средствами ОС.

Команда содержит одно целое число в диапазоне типа \texttt{int}. Числа 0 и 1 не являются ни простыми, ни составными: они пропускаются без завершения обработки. Некорректная строка отклоняется, после чего допускается ввод следующей команды. Конец стандартного ввода завершает обмен. Выходной файл создаётся либо обнуляется при запуске дочернего процесса.
